\section*{Appendix: Phase Prompt Templates}
\label{sec:appendix}

This appendix presents the phase prompt templates used by the Pigs API orchestration runtime. Each phase has a user-level instruction that is appended to the current user message at runtime. Both Chinese (zh) and English (en) versions are provided. The templates use \texttt{\{pre\_output\}} and \texttt{\{failure\_paths\}} as placeholder variables filled at runtime via string replacement.

\subsection*{A.1 Pre Phase (Planning) --- User Payload}

\paragraph*{Chinese (zh)}
\begin{lstlisting}[basicstyle=\ttfamily\scriptsize, frame=single]
以上是本任务要求
{failure_paths}
本次回答不需要具体执行任务，而是对任务进行执行前分析，
严格按要求输出这5个问题的回答
1）本任务需要哪些项目内部信息？
2）本任务需要哪些项目外部信息？
3）为本任务的执行制定执行计划
4）总结本任务的目标（goal），用于判断任务是否完成
5）判断本任务是否为简单问题。当且仅当1）-4）均不需要或
   极其简单时，才判定为简单问题，并必须给出理由，然后在
   回答最后单独一行输出 PIGEND。若不是简单问题，不要输出 PIGEND
\end{lstlisting}

\paragraph*{English (en)}
\begin{lstlisting}[basicstyle=\ttfamily\scriptsize, frame=single]
The text above is the task requirement.
{failure_paths}
Do not execute the task in this response. Analyze it before
execution and answer exactly these five questions:
1. What project-internal information is needed?
2. What external information is needed?
3. What is the execution plan for the task?
4. What is the task goal used to decide whether it is complete?
5. Decide whether this is a simple task. Only when questions 1-4
   need nothing or are extremely simple, classify it as simple,
   give the reason, and output PIGEND alone on the last line.
   If not simple, do not output PIGEND.
\end{lstlisting}

\subsection*{A.2 Executor Phase --- User Payload}

\paragraph*{Chinese (zh)}
\begin{lstlisting}[basicstyle=\ttfamily\scriptsize, frame=single]
以下是对本次任务在 1.内部信息的需求 2.外部信息的需求
3.执行计划 4.任务目标 5.任务难度方面的分析：
{pre_output}
完成对内部信息和外部信息的获取后，按照执行计划全力完成任务目标
\end{lstlisting}

\paragraph*{English (en)}
\begin{lstlisting}[basicstyle=\ttfamily\scriptsize, frame=single]
The following is the complete analysis of 1. internal information,
2. external information, 3. execution plan, 4. task goal, and
5. task difficulty:
{pre_output}
After obtaining the required internal and external information,
follow the execution plan and fully complete the task goal.
\end{lstlisting}

\subsection*{A.3 Post Phase (Review) --- User Payload}

\paragraph*{Chinese (zh)}
\begin{lstlisting}[basicstyle=\ttfamily\scriptsize, frame=single]
根据设定的目标，独立验收结果，对任务的执行情况进行全面审查。
先输出评审报告，进行如下判断：
1. 判定目标是否完全完成。仅当任务完成时，在回答最后输出 PIGEND
2. 判定任务是否是执行错误的，也就是当前离目标的距离比执行前
   更远，在回答最后输出 PIGFAIL
3. 如果是推进了任务，但没有完全完成任务；在输出评审报告后，
   继续执行任务，最后不输出标记
\end{lstlisting}

\paragraph*{English (en)}
\begin{lstlisting}[basicstyle=\ttfamily\scriptsize, frame=single]
Using the stated goal, independently verify the result and
comprehensively review execution of the task.
First output a review report, then make these judgments:
1. Decide whether the goal is fully complete. Only when it is
   complete, end the response with PIGEND.
2. Decide whether execution was wrong, meaning it moved farther
   from the goal than before execution. End with PIGFAIL.
3. If the task progressed but is not fully complete, continue
   executing it after the review report and do not output a marker.
\end{lstlisting}
